\subsection{Direction-aware ranking of non-translation residual absorbers}
\label{sec:q6-h-ranking-table-direction-aware-cross-layer}

\paragraph{Finite object.}
This record is a finite cross-layer ranking over the modeled-translation-complement residual of the $B^\ast_{Q6}$ protein-abundance component.  The classification is \emph{cross-layer relation}: it compares a nine-coordinate codon residual object with fetchable readouts such as measured ribosome-profiling translation efficiency, mRNA abundance, tRNA-copy tAI, protein half-life, localization labels, complex membership, transmembrane count, domain count, and post-translational-modification density.  For every organism and candidate readout $H$, the finite row stores the absorbed residual energy $A_H$, the coordinate-support cosine $C_H$, the raw direction-aware alignment $A_H C_H$, a stated description-length penalty, and a Bonferroni-adjusted null comparison when a matched null is available.  The scoring rule used here is
$$
\begin{aligned}
\operatorname{Score}(H) = A_H C_H - \lambda DL(H), \qquad \lambda = 0.01,
\end{aligned}
$$
with $DL(H)=\log_2(1+n_{\mathrm{columns}})$.  Thus one-column readouts have $DL=1$, whereas the twelve-column localization panel has $DL=3.700439718141092$.

\paragraph{Residual target and coordinate support.}
The residual carrier is built on the same nine coordinates in every computed organism: $K\_AAA$, $Arg\_AGR$, $Ile\_AUA$, $Leu\_CUN\_vs\_UUR$, $Leu\_UUA\_vs\_UUG$, $Ser\_UCR\_vs\_AGY$, $Ser\_UCA\_vs\_UCG$, $Thr\_ACR\_vs\_ACY$, and $f3\_stress$.  The complement computation reports full-rank control residualization with $24$ control columns and rank $23$ for each base organism.  The modeled-complement fractions are strongly organism-specific: $d=0.874628721049266$ in \emph{Danio rerio} on $14410$ joined rows, $d=0.6629960038578534$ in \emph{Escherichia coli} K-12 MG1655 on $3739$ rows, $d=0.9400341806302775$ in \emph{Gallus gallus} on $12418$ rows, $d=0.9208949355679455$ in \emph{Homo sapiens} on $19053$ rows, $d=0.8972672018746484$ in \emph{Mycobacterium smegmatis} on $4129$ rows, and $d=0.21173219572706714$ in \emph{Saccharomyces cerevisiae} on $5815$ rows.  The corresponding complement energies include $P_{Q\perp T\mathrm{mod}}^2=1005.3614101653926$ for danio, $287.00467177017873$ for \emph{E. coli}, $695.5324705841902$ for gallus, $1262.2237678771598$ for human, $164.9629139386273$ for mycobacterium, and $186.88597748267557$ for yeast.

\paragraph{Top rows by raw direction-aware alignment.}
The ranking does not find a single shared absorber.  In danio, measured TE is the top row by $A_HC_H$, with $A_H=0.017594290015070787$, $C_H=0.6573753096184687$, $A_HC_H=0.011566051846174291$, and $\operatorname{Score}=0.001566051846174291$ on $13140$ rows; mRNA follows with $A_HC_H=0.01073458857797759$ and $\operatorname{Score}=0.000734588577977589$.  Both exceed their Bonferroni thresholds, respectively $0.000719994971498518$ and $0.000587830414454545$.  In human, the best one-column row is complex membership, with $A_H=0.007527199987780841$, $C_H=0.8723071163246565$, $A_HC_H=0.006566030115340095$, Bonferroni threshold $0.0002646171322365322$, and $\operatorname{Score}=-0.0034339698846599055$ on $18335$ rows.  Human mRNA is close, with $A_HC_H=0.006224777217900546$ on $11554$ rows.  Yeast is much weaker: mRNA has $A_H=0.001897605472160837$, $C_H=0.6040571890503664$, $A_HC_H=0.0011462622274400686$, and $\operatorname{Score}=-0.008853737772559931$ on $5128$ rows.  In \emph{E. coli}, mRNA has $A_HC_H=0.0015297976209788734$ but does not exceed its Bonferroni threshold $0.001860998420625413$; measured TE is still lower, with $A_HC_H=0.0005536661857560242$.  Gallus has only feature-panel rows in this table: complex membership is the raw leader with $A_HC_H=0.004241514781443171$, while PTM density has larger absorbed energy $A_H=0.023360446890988432$ but a low direction cosine $C_H=0.10469593979758543$.

\paragraph{Direction filter.}
The cosine term is not decorative; it changes the reading of high-energy rows.  Human PTM density has the largest human one-column absorbed energy, $A_H=0.014884825920744861$, and exceeds a Bonferroni threshold of $0.0004629487556325413$, but its direction cosine is only $C_H=0.3604440741749213$, giving $A_HC_H=0.0053651472982577525$ and placing it behind complex membership and mRNA in raw direction-aware order.  The same issue is sharper for human measured TE: $A_H=0.002894714925796667$ exceeds its Bonferroni threshold $0.0004670613303535708$, yet $C_H=-0.10813537216189718$, so $A_HC_H=-0.00031302107580362116$ and the Score is $-0.01031302107580362$.  Human measured tRNA-copy tAI is also wrong-direction, with $A_H=5.9525256039488626\times 10^{-5}$ and $C_H=-0.3998537725071249$.  Yeast domain count is similarly negative-direction, with $A_H=7.029698692511865\times 10^{-5}$ and $C_H=-0.02003766381646423$.

\paragraph{Penalty and null behavior.}
The description-length term separates finite absorption from compact explanatory value.  Under the stated penalty, only the two danio one-column rows have positive Scores: measured TE at $0.001566051846174291$ and mRNA at $0.000734588577977589$.  The top human one-column rows remain negative despite strong Bonferroni excess, and yeast remains negative throughout.  The twelve-column localization readout is especially illustrative.  Human localization has $A_H=0.009254913329269913$, $C_H=0.7450003925723202$, $A_HC_H=0.006894914063528885$, exceeds a Bonferroni threshold of $0.002033625458961835$, and uses $16008$ joined rows; after $DL=3.700439718141092$, its Score is $-0.030109483117882037$.  Yeast localization similarly exceeds its threshold, $0.006295217548252396$, with $A_H=0.011225105608516642$ and $C_H=0.4247910216132792$, but scores $-0.03223607310225219$.  Danio localization does not exceed its threshold: $A_H=0.009508864544725801$ against Bonferroni threshold $0.013392249164674856$ on $1984$ joined rows.

\paragraph{Readout provenance and row loss.}
The measured TE and mRNA rows use organism-specific joins.  For danio, both measured TE and mRNA use $13140$ rows rather than the $14410$ base CDS-protein rows; for human, both use $11554$ rows rather than $19053$; for yeast, both use $5128$ rows rather than $5815$; for \emph{E. coli}, both use $3415$ rows rather than $3739$.  Measured tRNA-copy tAI keeps the full complement join in the three organisms where it is computed: $14410$ danio rows, $19053$ human rows, and $5815$ yeast rows, with anticodon table sizes $57$, $47$, and $41$ and tRNA-gene counts $8846$, $431$, and $273$.  Turnover is narrower.  Yeast turnover joins $3640$ rows, with a reported turnover join hit rate $0.8818572656921754$ and CDS join hit rate $0.9156038419146592$; its absorption is $0.0009086031889974897$ in the ranking table and its independent turnover run reports $d=0.2232704509412724$.  Human turnover joins $8057$ rows out of a $19053$-row complement reference, with CDS hit rate $0.9779294769799313$, but the dof-matched null does not support excess: absorbed turnover is $3.5740764900116538\times 10^{-6}$, null mean is $0.0001262572958089998$, and null $95$th percentile is $0.0004342285066808464$.

\paragraph{Cross-organism comparison.}
The per-organism best absorber is not stable.  The table records $H_o^\ast=\mathrm{measured\_te}$ for danio, $H_o^\ast=\mathrm{complex\_member}$ for human and gallus, $H_o^\ast=\mathrm{mrna}$ for yeast, and no selected best absorber for \emph{E. coli} or mycobacterium.  This agrees with the broader organism-specificity checks: the dominant-coordinate contingency over domain has $n=12$, $\chi^2=9.399999999999999$, Cramer's $V=0.6258327785172862$, but permutation $p=0.5814185814185814$, so the observed heterogeneity is not resolved as a domain law.  Continuous meta-tests are likewise descriptive rather than law-like: $M_{QTP}$ fraction versus GC3 has $r=-0.3062616500469447$ and two-sided permutation $p=0.35264735264735264$; versus log proteome size has $r=-0.033628115789504955$ and $p=0.913086913086913$; versus log tRNA-gene count has $r=0.06029405206037823$ and $p=0.8621378621378621$.  A separate optimal-codon direction survey reports $15$ computed organisms out of $16$ requested, $41$ synonymous design columns, and mean pairwise cosine $0.2629282074265475$, classified as clustered rather than universal.

\paragraph{Related finite checks.}
The measured mediation table gives another view of why this ranking should stay descriptive.  Measured TE mediation is large in yeast, with $M_{QTP,\mathrm{meas}}=0.4793069656345167$, but small in \emph{E. coli} at $0.04091158340950673$, human at $0.0027966485036786744$, and clipped to $0$ in danio after a raw value $-0.04801486048721136$.  The modeled-complement residual remains large in several organisms even after translation-side summaries.  Ordered decompositions report unexplained fractions $0.9756319042010854$ in danio, $0.9942823785118229$ in \emph{E. coli}, $0.9736442326306781$ in human, and $0.8676586345294415$ in yeast.  The minimal-compression test has median compressibility $0.026$ across computed organisms, with organism values $0.024368095798914635$, $0.005717621488177116$, $0.02635576736932188$, and $0.13234136547055853$; all listed permutation values are $0.000999000999000999$, but the conclusion is low compression by the measured low-dimensional readout set, not a compact common mechanism.  The human tissue-structure check also stays weakly descriptive: across $13$ adult tissue proteome readouts, the within-class minus between-class cosine difference is $0.007926963193379821$ with one-sided permutation $p=0.3006993006993007$.  For the yeast $f3\_stress$ coordinate, the best independent-readout subset reaches coverage $0.22618445168998694$ against a selection-aware null $95$th percentile $0.02697141300639046$, but $0.7738155483100131$ remains unexplained after the full subset.

\paragraph{Interpretation boundary.}
The only asserted output of this packet is the finite ranking readback: availability, $A_H$, $C_H$, $A_HC_H$, Score, Bonferroni status, rank, and per-organism best absorber under the declared rows and penalties.  A positive Bonferroni flag is not a mechanism, a high direction cosine is not a biological pathway, and a negative Score is not evidence that the corresponding biological layer is absent.  Missing rows remain missing, not zero and not within-null.  No translation realization, folding statement, physical-admissibility claim, molecular-function claim, phenotype claim, causal pathway, universal absorber, or global biological law follows from these numbers unless a separate matched measurement and statistic supplies that stronger claim.  The present result is therefore an honest negative-to-weak finite table: some organism-specific readouts absorb small pieces of the modeled-translation-complement residual, the direction filter removes several high-energy wrong-direction rows, the description-length penalty leaves almost no compact positive explanations, and no shared universal $H$ is found.

\origin{ai}
